\documentclass[12pt,a4paper]{extarticle}

\usepackage[utf8]{inputenc}
\usepackage[german]{babel}
\usepackage{geometry}
\usepackage{setspace}
\usepackage{fancyhdr}
\usepackage[parfill]{parskip} 

\geometry{a4paper, margin=2.5cm}
\pagestyle{fancy}
\fancyhf{}
\lhead{Chat-Co}
\chead{Version 3.0}
\rhead{Lastenheft}
\rfoot{\thepage}
\lfoot{Team: Alexander Rutz / Maximilian Paesold / Sebastian Weigl}
\renewcommand{\headrulewidth}{0.4pt}
\renewcommand{\footrulewidth}{0.4pt}

\onehalfspacing

\title{
    \textbf{Lastenheft Chat-Co} \\
    \large Sichere On-Premise Kommunikationslösung
}
\author{Alexander Rutz (PM), Maximilian Paesold, Sebastian Weigl}
\date{Version 3.0 \\ Status: Überarbeitet}

\begin{document}
\maketitle

\begin{center}
  \begin{tabular}{|l|l|l|l|l|}
    \hline
    Version & Autor & QS & Datum & Kommentar \\
    \hline
    1.0 & Projektteam & Rutz & 12.02.2026 & Erster Entwurf \\
    \hline
    1.1 & Projektteam & Rutz & 15.02.2026 & Überarbeitung Struktur \\
    \hline
    2.0 & Projektteam & Rutz & 20.02.2026 & Anpassung gemäß Team Feedback \\
    \hline
    3.0 & Projektteam & Rutz & 27.02.2026 & Anpassungen gemäß Feedback\\
    \hline
\end{tabular}
\end{center}

\newpage
\tableofcontents
\newpage


\section{Einführung}

Digitale Kommunikation ist ein zentraler Bestandteil moderner Unternehmensprozesse. Viele Organisationen nutzen Cloud-basierte Lösungen, wodurch sensible Daten außerhalb der eigenen Infrastruktur gespeichert werden.
Das Projekt Chat-Co verfolgt daher einen konsequenten On-Premise-Ansatz. Sämtliche Kommunikations- und Organisationsdaten verbleiben innerhalb der unternehmenseigenen IT-Infrastruktur. Dadurch wird die vollständige Datenhoheit gewährleistet und bestehende Sicherheitsrichtlinien können eingehalten werden.

\vspace{1cm}
\section{Zielbestimmung}

Ziel des Projektes ist die Entwicklung einer sicheren, unternehmensinternen Kommunikationsplattform.
Das System soll eine strukturierte und datenschutzkonforme Kommunikation innerhalb einer Organisation ermöglichen. Dabei dürfen sämtliche Daten ausschließlich innerhalb der eigenen IT-Infrastruktur verarbeitet werden.
Darüber hinaus soll die interne Zusammenarbeit durch Echtzeitkommunikation sowie Termin- und Ressourcenverwaltung effizienter gestaltet werden.

\vspace{1cm}
\section{Produkteinsatz}

Chat-Co wird in Unternehmen, Behörden und Bildungseinrichtungen eingesetzt, die erhöhte Anforderungen an Datenschutz und Datensicherheit haben. Das Produkt dient der internen Kommunikation von Mitarbeitern sowie der Organisation von Meetings und Ressourcen. Zielgruppe sind Administratoren, Mitarbeiter und Projektteams innerhalb einer Organisation.

\newpage
\section{Produktübersicht}

Chat-Co ist eine On-Premise Kommunikationsplattform zur Verwaltung von Textnachrichten, Dateien, Terminen und Ressourcen.
Das System wird innerhalb der bestehenden IT-Infrastruktur betrieben und integriert vorhandene Benutzerverzeichnisse.

\section{Produktfunktionen}

\subsection{Benutzer- und Sicherheitsfunktionen}

/LF10/ \textbf{LDAP-Anbindung:} Das System muss die Anmeldung eines Benutzers über LDAP ermöglichen.

/LF20/ \textbf{Rechteverwaltung bereitstellen:} Das System muss eine rollenbasierte Rechteverwaltung für Administratoren, Mitarbeiter und Gäste bereitstellen.

/LF30/ \textbf{Benutzerprofile verwalten:} Das System muss Benutzerprofile mit Namen, Statusanzeige und optionalem Avatar verwalten.

/LF40/ \textbf{Inaktivitäts-Abmeldung erzwingen:} Das System muss Benutzer bei Inaktivität automatisch abmelden.

\subsection{Kommunikationsfunktionen}

/LF50/ \textbf{Echtzeit-Kommunikation ermöglichen:} Das System muss eine Echtzeit-\\Textkommunikation zwischen Benutzern ermöglichen.

/LF60/ \textbf{Nachrichtenbearbeitung unterstützen:} Das System muss das Bearbeiten und Löschen gesendeter Nachrichten unterstützen.

/LF70/ \textbf{Textformatierung integrieren:} Das System muss die Verwendung von Emojis und einfacher Textformatierung unterstützen.

\subsection{Datei- und Kanalverwaltung}

/LF80/ \textbf{Kanalverwaltung ermöglichen:} Das System muss das Erstellen und Verwalten von Kommunikationskanälen ermöglichen.

/LF90/ \textbf{Kanalzuordnung durchführen:} Das System muss Benutzer einzelnen Kanälen zuordnen können.

/LF100/ \textbf{Dateiaustausch absichern:} Das System muss den sicheren Austausch von Dateien innerhalb von Chats ermöglichen.

\newpage
\subsection{Meeting- und Organisationsfunktionen}

/LF110/ \textbf{Meeting-Erstellung ermöglichen:} Das System muss das Erstellen von Meeting-Anfragen ermöglichen.

/LF120/ \textbf{Einladungsversand steuern:} Das System muss Einladungen an ausgewählte Benutzer versenden können.

/LF130/ \textbf{Raumbuchung verwalten:} Das System muss verfügbare Räume anzeigen und Buchungen verwalten.

/LF140/ \textbf{Doppelbuchungs-Vermeidung garantieren:} Das System muss Doppelbuchungen von Räumen verhindern.  
 

\vspace{1cm}
\section{Produktdaten}

/LD10/ Das System speichert Benutzerdaten wie Benutzernamen, LDAP-UIDs, E-Mail-Adressen, Anzeigenamen und den Aktivitätsstatus.

/LD20/ Es werden Rollen und Berechtigungen sowie die m:n-Zuordnungen zwischen Benutzern und Rollen gespeichert.

/LD30/ Die Speicherung von Chatdaten umfasst Konversationen, Teilnehmerlisten und Nachrichten inklusive Metadaten wie Zeitstempel, Typ, Löschstatus und Referenzen.

/LD40/ Dateien und Anhänge werden zusammen mit Metadaten (Name, Pfad, Typ, Größe, Upload-Zeit) und der Zuordnung zu Benutzern oder Nachrichten hinterlegt.

/LD50/ Das System speichert Kalender- und Meetingdaten inklusive Titel, Zeiträumen, Orten/Links sowie Teilnehmern und deren Status siehe LD10.

/LD60/ Es werden Daten zu Räumen und Ressourcen gespeichert, um Buchungen und Belegungen je Zeitraum abzubilden.

/LD70/ Systemprotokolle werden mit Zeitstempel, Log-Level, Aktion und optionalem Bezug zu Benutzern oder Objekten erfasst.  

\newpage
\section{Produktleistungen}

/LL10/ Die Anmeldung eines Benutzers darf eine Interaktionszeit von fünf Sekunden nicht überschreiten.  

/LL20/ Chatnachrichten müssen innerhalb von fünf Sekunden beim Empfänger angezeigt werden.  

/LL30/ Das System muss mindestens 50 gleichzeitige Benutzer unterstützen.  

/LL40/ Es müssen mindestens 10.000 Nachrichten pro Monat verarbeitet werden können.  

/LL50/ Alle Datenübertragungen müssen verschlüsselt erfolgen.  

/LL60/ Das System muss regelmäßige Datensicherungen ermöglichen.  
\vspace{0.5cm}
\section{Qualitätsanforderungen}

\begin{center}
  \begin{tabular}{|l|c|c|c|c|}
\hline
Produktqualität & Sehr gut & Gut & Normal & Irrelevant \\ \hline
Funktionalität & X & & & \\ \hline
Zuverlässigkeit & X & & & \\ \hline
Benutzbarkeit & & X & & \\ \hline
Effizienz & & X & & \\ \hline
Änderbarkeit & & & X & \\ \hline
Übertragbarkeit & & X & & \\ \hline
\end{tabular}  
\end{center}

\subsection*{Erläuterungen}

\textbf{Funktionalität (Sehr gut):}  
Alle Kernfunktionen müssen zuverlässig und vollständig verfügbar sein.

\textbf{Zuverlässigkeit (Sehr gut):}  
Das System muss stabil laufen und Daten konsistent speichern.

\textbf{Benutzbarkeit (Gut):}  
Die Oberfläche soll intuitiv bedienbar sein und schnelle Einarbeitung ermöglichen.

\textbf{Effizienz (Gut):}  
Das System soll ressourcenschonend arbeiten und eine maximale Reaktionszeiten von 5 Sekunden haben.

\textbf{Änderbarkeit (Normal):}  
Die modulare Architektur ermöglicht Erweiterungen ohne Systemgefährdung.

\textbf{Übertragbarkeit (Gut):}  
Die Java-Basis ermöglicht plattformübergreifenden Betrieb.


\newpage
\section{Glossar}

Hier werden die wichtigsten Fachbegriffe für das Projekt Chat-Co erläutert:

\begin{description}
\item[LDAP] \textit{Lightweight Directory Access Protocol}. Ein Standard-Protokoll zur Abfrage von Nutzerinformationen aus einem Verzeichnisdienst (z. B. Active Directory).
\item[On-Premise] Software, die auf eigenen Servern im eigenen Rechenzentrum betrieben wird, statt sie als Dienst aus einer externen Cloud zu beziehen.
\item[Spring Boot] Ein Java-basiertes Framework, das die Erstellung von Microservices und Web-Backends durch Konventionen vereinfacht.
\item[Vaadin] Ein Web-Framework für Java, das es ermöglicht, moderne Benutzeroberflächen vollständig in Java zu programmieren.
\item[WebSockets] Ein Kommunikationsprotokoll, das eine dauerhafte, bidirektionale Verbindung zwischen Client und Server ermöglicht.
\item[MVC] \textit{Model-View-Controller}. Ein Architekturmuster zur Trennung von Datenmodell, Benutzeroberfläche und Steuerungslogik.
\item[PostgreSQL] Ein objektrelationales Datenbankmanagementsystem.
\item[TLS] \textit{Transport Layer Security}. Ein Verschlüsselungsprotokoll zur sicheren Datenübertragung.
\item[JAR-Datei] \textit{Java Archive}. Ein Dateiformat zur Bündelung von Java-Anwendungen.
\end{description}

\end{document}